\centerline{\bf \Large  РЕШЕНИЯ и КРИТЕРИИ}

\centerline{\bf \Large  10 КЛАСС}

\problem{\textbf{1}} 
Ответ: 2021. Решение. Сумма равна $(10-7)+(100-7)+\ldots+\left(10^{2026}-7\right)=11 \ldots 110-$ $14182=11 \ldots 11096928$, где в числе $11 \ldots 110-2026$ единиц, и, следовательно, в числе $11 \ldots 11096928-2021$ единиц.

\textbf{Критерии.} \\
\textit{7 баллов} $-$ Полностью верное решение\\
\textit{6 баллов} $-$ Полностью верное решение, но в конце допущена арифметическая ошибка\\
\textit{3 балла} $-$ Выписано выражение $(10-7)+(100-7)+\ldots+\left(10^{2026}-7\right)$ без дальнейшего решения\\
\textit{0 баллов} $-$  Только ответ\\ 


\problem{\textbf{2}} Так как все треугольники равнобедренные и угол подъёма одинаков, это означает, что углы при основаниях этих равнобедренных треугольников равны. Значит, все равнобедренные треугольники подобны. Так как каждая последующая высота в 5 раз больше, то и сумма боковых сторон в 5 раз больше, соответственно, времени на переход через каждый последующий холм требуется в 5 раз больше. Подъём занимает 2 часа, а спуск - в два раза меньше. Значит, спускается за 1 час. Таким образом, первый холм перейдет за 3 часа, второй - за $3\cdot5$ часа, третий - за $3\cdot5^2$ часов, ... , $n$-ый - за $3\cdot5^n$ часов. Вся эта сумма не должна превышать $24\cdot8000$ часов.

Решая неравенство через сумму геометрической прогрессии получаем, что $n=6$\par

\textbf{Критерии.} \\
\textit{0-4 баллов} $-$ Решение пункта А\\
\textit{0-1 балла} $-$ Составление условия пункта Б \\
\textit{0-2 баллов} $-$ Решение пункта Б\\
\textit{0 баллов} $-$ Отсутствие решения пункта А (сразу НОЛЬ! баллов)\\
\textit{0 баллов} $-$ Любой ответ без объяснений\\

\problem{\textbf{3}}
$X=\sqrt[3]{4}+\sqrt[3]{2}+\sqrt[3]{1} \\ X(\sqrt[3]{2}-1)=(\sqrt[3]{2}-1) \sqrt[3]{4}+\sqrt[3]{2}+\sqrt[3]{1} \\ X=\frac{1}{\sqrt[3]{2}-1}$

Преобразуем требуемое выражение 

$$\dfrac{3X^2 + 3X + 1}{X^3} = \dfrac{(X+1)^3-X^3}{X^3} = \left(\dfrac{X+1}{X}\right)^3 - 1 = 2 - 1$$

\textbf{Критерии.} \\
\textit{7 баллов} $-$ Получен ответ полным раскрытием скобок\\
\textit{2 балла} $-$ Получено равенство  $X=\frac{1}{\sqrt[3]{2}-1}$\\
\textit{1 балл} $-$ Получено равенство $\dfrac{3}{x^2} + \dfrac{3}{x} + \dfrac{1}{x^3}=\left(\dfrac{X+1}{X}\right)^3$\\ 
\textit{0 баллов} $-$ Только ответ\\


\problem{\textbf{4}} С одной стороны, рыбки каждого цвета должны встречаться по крайней мере в
трёх аквариумах из шести, поскольку если рыбки какого-то цвета лежат не более чем в двух
аквариумах, то в оставшихся аквариумах, которых — не менее четырёх, рыбок этого цвета нет.
Поэтому всего рыбок должно быть не менее $3\times26$.
С другой стороны, если взять по 3 рыбки каждого из 26 цветов и распределить их так,
чтобы никакие рыбки одного цвета не были в одном аквариуме, то рыбка любого цвета
найдётся в любых четырёх аквариумах (так как остальных аквариумов всего два). Нужное расположение
рыбок получится, например, если всех рыбок разбить на 3 группы по 26 рыбок
всех цветов и каждую группу распределить поровну (по 13 рыбок) в два аквариума. \par

\textbf{Критерии.} \\
\textit{7 баллов} $-$ Полностью верное решение\\
\textit{4 баллов} $-$ Верный пример по 13\\
\textit{3 баллов} $-$ Верная оценка, не менее $3\times26$\\
\textit{0-1 балл} $-$ За все остальные попытки решить задачу\\
\textit{0 баллов} $-$ Любой ответ без объяснений\\


\problem{\textbf{5}} 
Решение. Обозначим угол $A$ через $\alpha$. Пусть перпендикуляр из точки $A$ к $A C$ пересекает $A B$ в точке $Q$, а $R$ - такая точка на $A B$, что $R B=C B$. Тогда, как нетрудно видеть, $\angle C R B=\angle A Q C=90^{\circ}-\alpha$, поэтому $C Q=C R$. Поскольку $C R=2 B C \sin \alpha$, а $C Q=A Q \sin \alpha$, получаем $A Q=2 B C=B P$. Таким образом, $M-$ середина отрезка $P Q$, а перпендикуляр из $M$ к $A C$ - средняя линия треугольника $C P O$, из чего и следует утверждение задачи.

\textbf{Критерии.} \\
\textit{7 баллов} $-$ Полностью верное решение\\
\textit{3 балла} $-$ Проведена прямая из вершины $C$ параллельная высоте из $M$ к стороне $AC$\\
\textit{1-3 балла} $-$  Доп построения/подсчёт углов, которые в теории могу привести к решению, но оно не предъявлено.\\ 
\textit{0-1 балл} $-$ За все остальные попытки решить задачу\\

